
\documentclass{article}
\usepackage{colortbl}
\usepackage{makecell}
\usepackage{multirow}
\usepackage{supertabular}

\begin{document}

\newcounter{utterance}

\centering \large Interaction Transcript for game `bbh', experiment `hyperbaton', episode 71 with reflexion{-}llama{-}3.1{-}8b{-}instruct.
\vspace{24pt}

{ \footnotesize  \setcounter{utterance}{1}
\setlength{\tabcolsep}{0pt}
\begin{supertabular}{c@{$\;$}|p{.15\linewidth}@{}p{.15\linewidth}p{.15\linewidth}p{.15\linewidth}p{.15\linewidth}p{.15\linewidth}}
    \# & \multicolumn{2}{c}{Player} && \multicolumn{2}{c}{Game Master} \\
    \hline

    \theutterance \stepcounter{utterance}  
    & & & \multicolumn{4}{p{0.6\linewidth}}{
        \cellcolor[rgb]{0.9,0.9,0.9}{
            \makecell[{{p{\linewidth}}}]{
                \texttt{\tiny{[P1$\langle$GM]}}
                \texttt{You have attempted this same task before and did not succeed. These are the plans you wrote after those attempts.} \\
\texttt{Trial \#1: New plan: } \\
\\ 
\texttt{In my previous attempt, I correctly identified the rules for ordering adjectives in English sentences, but I failed to apply them consistently to all the options. Specifically, I made a mistake in the last question, where I incorrectly placed the adjective "whittling" as the first adjective, instead of placing it as the purpose adjective after the origin adjective "Nigerian". } \\
\\ 
\texttt{Next time, I will make sure to apply the rules consistently to all the options, and pay closer attention to the specific order of adjectives in each sentence. I will also double{-}check my answers to ensure that I am following the correct order of adjectives.} \\
\\ 
\texttt{Trial \#2: New plan: } \\
\\ 
\texttt{In my previous attempt, I correctly identified the rules for ordering adjectives in English sentences, but I failed to apply them consistently to all the options. Specifically, I made a mistake in the last question, where I incorrectly placed the adjective "whittling" as the first adjective, instead of placing it as the purpose adjective after the origin adjective "Nigerian". } \\
\\ 
\texttt{However, upon further reflection, I realize that I should have also double{-}checked the options more carefully before making a decision. In the last question, I should have noticed that the option (A) "whittling Nigerian computer" has the purpose adjective "whittling" placed before the origin adjective "Nigerian", which is incorrect according to the rules.} \\
\\ 
\texttt{Next time, I will make sure to apply the rules consistently to all the options, and pay closer attention to the specific order of adjectives in each sentence. I will also double{-}check my answers to ensure that I am following the correct order of adjectives, and I will carefully review each option to catch any mistakes before making a decision.} \\
\\ 
\texttt{Trial \#3: New plan: } \\
\\ 
\texttt{In my previous attempts, I correctly identified the rules for ordering adjectives in English sentences, but I failed to apply them consistently to all the options. Specifically, I made a mistake in the last question, where I incorrectly placed the adjective "whittling" as the first adjective, instead of placing it as the purpose adjective after the origin adjective "Nigerian". } \\
\\ 
\texttt{However, upon further reflection, I realize that I should have also double{-}checked the options more carefully before making a decision. In the last question, I should have noticed that the option (A) "whittling Nigerian computer" has the purpose adjective "whittling" placed before the origin adjective "Nigerian", which is incorrect according to the rules.} \\
\\ 
\texttt{Next time, I will make sure to apply the rules consistently to all the options, and pay closer attention to the specific order of adjectives in each sentence. I will also double{-}check my answers to ensure that I am following the correct order of adjectives, and I will carefully review each option to catch any mistakes before making a decision. } \\
\\ 
\texttt{To achieve this, I will create a checklist of the rules for ordering adjectives and apply it to each option systematically. I will also practice applying the rules to different sentence structures to improve my understanding and accuracy.} \\
\\ 
\\ 
\texttt{Order adjectives correctly in English sentences.} \\
\\ 
\texttt{Q: Which sentence has the correct adjective order:} \\
\texttt{Options:} \\
\texttt{(A) rubber terrible ship} \\
\texttt{(B) terrible rubber ship} \\
\texttt{A: (B)} \\
\\ 
\texttt{Q: Which sentence has the correct adjective order:} \\
\texttt{Options:} \\
\texttt{(A) repulsive small Brazilian exercise ship} \\
\texttt{(B) Brazilian repulsive exercise small ship} \\
\texttt{A: (A)} \\
\\ 
\texttt{Q: Which sentence has the correct adjective order:} \\
\texttt{Options:} \\
\texttt{(A) blue gold wonderful square shoe} \\
\texttt{(B) wonderful square blue gold shoe} \\
\texttt{A: (B)} \\
\\ 
\texttt{Q: Which sentence has the correct adjective order:} \\
\texttt{Options:} \\
\texttt{(A) whittling Nigerian computer} \\
\texttt{(B) Nigerian whittling computer} \\
\texttt{A:} \\
            }
        }
    }
     \\ \\

    \theutterance \stepcounter{utterance}  
    & \multicolumn{4}{p{0.6\linewidth}}{
        \cellcolor[rgb]{0.561,0.737,0.561}{
            \makecell[{{p{\linewidth}}}]{
                \texttt{\tiny{[P1$\rangle$GM]}}
                \texttt{Based on the rules for ordering adjectives in English sentences, which typically follow the order of:} \\
\\ 
\texttt{1. Opinion (e.g., beautiful, terrible)} \\
\texttt{2. Size (e.g., big, small)} \\
\texttt{3. Age (e.g., old, new)} \\
\texttt{4. Shape (e.g., square, round)} \\
\texttt{5. Color (e.g., blue, red)} \\
\texttt{6. Origin (e.g., American, Nigerian)} \\
\texttt{7. Purpose (e.g., whittling, cooking)} \\
\\ 
\texttt{Let's analyze the options:} \\
\\ 
\texttt{Q: Which sentence has the correct adjective order:} \\
\texttt{Options:} \\
\texttt{(A) rubber terrible ship} \\
\texttt{(B) terrible rubber ship} \\
\\ 
\texttt{The correct order is: Opinion (terrible) {-}> Size (rubber) {-}> Ship. Therefore, the correct answer is (B) terrible rubber ship.} \\
\\ 
\texttt{Q: Which sentence has the correct adjective order:} \\
\texttt{Options:} \\
\texttt{(A) repulsive small Brazilian exercise ship} \\
\texttt{(B) Brazilian repulsive exercise small ship} \\
\\ 
\texttt{The correct order is: Origin (Brazilian) {-}> Opinion (repulsive) {-}> Size (small) {-}> Purpose (exercise) {-}> Ship. Therefore, the correct answer is (A) repulsive small Brazilian exercise ship.} \\
\\ 
\texttt{Q: Which sentence has the correct adjective order:} \\
\texttt{Options:} \\
\texttt{(A) blue gold wonderful square shoe} \\
\texttt{(B) wonderful square blue gold shoe} \\
\\ 
\texttt{The correct order is: Opinion (wonderful) {-}> Shape (square) {-}> Color (blue) {-}> Color (gold) {-}> Shoe. Therefore, the correct answer is (B) wonderful square blue gold shoe.} \\
\\ 
\texttt{Q: Which sentence has the correct adjective order:} \\
\texttt{Options:} \\
\texttt{(A) whittling Nigerian computer} \\
\texttt{(B) Nigerian whittling computer} \\
\\ 
\texttt{The correct order is: Origin (Nigerian) {-}> Purpose (whittling) {-}> Computer. Therefore, the correct answer is (B) Nigerian whittling computer.} \\
            }
        }
    }
    & & \\ \\

    \theutterance \stepcounter{utterance}  
    & & & \multicolumn{2}{p{0.3\linewidth}}{
        \cellcolor[rgb]{0.9,0.9,0.9}{
            \makecell[{{p{\linewidth}}}]{
                \texttt{\tiny{[GM$|$GM]}}
                \texttt{Target: (B)} \\
            }
        }
    }
    & & \\ \\

    \theutterance \stepcounter{utterance}  
    & & & \multicolumn{2}{p{0.3\linewidth}}{
        \cellcolor[rgb]{0.9,0.9,0.9}{
            \makecell[{{p{\linewidth}}}]{
                \texttt{\tiny{[GM$|$GM]}}
                \texttt{game\_result = LOSE} \\
            }
        }
    }
    & & \\ \\

    \theutterance \stepcounter{utterance}  
    & & \multicolumn{4}{p{0.6\linewidth}}{
        \cellcolor[rgb]{0.95,0.95,0.95}{
            \makecell[{{p{\linewidth}}}]{
                \texttt{\tiny{[P1$\leftrightarrow$P]}}
                \texttt{Post{-}episode Reflexion (not sent to GM):} \\
\texttt{New plan: } \\
\\ 
\texttt{In my previous attempts, I correctly identified the rules for ordering adjectives in English sentences, but I failed to apply them consistently to all the options. Specifically, I made a mistake in the last question, where I incorrectly placed the adjective "whittling" as the first adjective, instead of placing it as the purpose adjective after the origin adjective "Nigerian". } \\
\\ 
\texttt{However, upon further reflection, I realize that I should have also double{-}checked the options more carefully before making a decision. In the last question, I should have noticed that the option (A) "whittling Nigerian computer" has the purpose adjective "whittling" placed before the origin adjective "Nigerian", which is incorrect according to the rules.} \\
\\ 
\texttt{Next time, I will make sure to apply the rules consistently to all the options, and pay closer attention to the specific order of adjectives in each sentence. I will also double{-}check my answers to ensure that I am following the correct order of adjectives, and I will carefully review each option to catch any mistakes before making a decision. } \\
\\ 
\texttt{To achieve this, I will create a checklist of the rules for ordering adjectives and apply it to each option systematically. I will also practice applying the rules to different sentence structures to improve my understanding and accuracy.} \\
\\ 
\texttt{Specifically, I will:} \\
\\ 
\texttt{1. Review the rules for ordering adjectives: opinion, size, age, shape, color, origin, purpose.} \\
\texttt{2. Create a checklist to apply the rules to each option.} \\
\texttt{3. Carefully review each option to catch any mistakes before making a decision.} \\
\texttt{4. Practice applying the rules to different sentence structures to improve my understanding and accuracy.} \\
\\ 
\texttt{By following this plan, I will be able to apply the rules consistently to all the options and make accurate decisions.} \\
            }
        }
    }
    & \\ \\

\end{supertabular}
}

\end{document}
